% ============================================================================
% Fichier  : partie2/P02_C04_droit_preuve.tex : FORN (LPI342)
% Sources : loi 2010/012, loi 2024/017 ; Malabo 2014 ; Budapest 2001.
% ============================================================================

\chapter{Le Droit de la Preuve Numérique}
\label{chap:droit}
\chaptermark{Le droit de la preuve numérique}

\epigraph{%
    « Une preuve obtenue hors du droit\\
    affaiblit la cause qu'elle prétend servir. »%
}{Principe du cours}

\niveauChapitre{Fondamental}{Chapitre~\ref{chap:ethique}~; notions de droit.}

\begin{boxobjectifs}
\begin{itemize}
    \item Connaître le cadre camerounais de la preuve numérique.
    \item Situer la loi 2010/012 et la loi 2024/017.
    \item Articuler avec Malabo et Budapest.
    \item Comprendre le rôle du juge et de l'expert.
    \item Déduire les règles à respecter pendant l'enquête.
\end{itemize}
\end{boxobjectifs}

\bigskip

% ============================================================================
\section*{Question directrice : que faut-il respecter pour que la preuve tienne devant le juge ?}
% ============================================================================

L'enquête se déroule sous le regard du droit~: le titre pour agir, le respect
des données personnelles, les règles de la preuve. Ce chapitre pose ce cadre.

\medskip
\textbf{Plan du chapitre.}
\begin{enumerate}
    \item Le cadre camerounais (section~\ref{sec:c04-cameroun}).
    \item Les données personnelles (section~\ref{sec:c04-donnees}).
    \item Malabo et Budapest (section~\ref{sec:c04-international}).
    \item Le juge et l'expert (section~\ref{sec:c04-juge}).
\end{enumerate}

\separateur

% ============================================================================
\section{Le Cadre Camerounais}
\label{sec:c04-cameroun}
% ============================================================================

\begin{boxjuris}
La \loicam{loi n°~2010/012 du 21~décembre~2010} relative à la cybersécurité
et à la cybercriminalité réprime l'accès et le maintien frauduleux dans un
système et encadre la preuve électronique. L'accès à un support sans titre
(réquisition, mandat, autorisation) est illicite~; la preuve ainsi obtenue
risque l'irrecevabilité.
\end{boxjuris}

% ============================================================================
\section{Les Données Personnelles}
\label{sec:c04-donnees}
% ============================================================================

\begin{boxjuris}
La \loicam{loi n°~2024/017 du 23~décembre~2024} régit la protection des
données à caractère personnel. Un support d'enquête contient des données
personnelles~: on les traite dans le périmètre de la réquisition, on ne
diffuse rien au-delà, et l'on protège les tiers.
\end{boxjuris}

% ============================================================================
\section{Malabo et Budapest}
\label{sec:c04-international}
% ============================================================================

\begin{tableaucours}{L{3.6cm}L{8.6cm}}{%
    \tableauHeader{Instrument} &
    \tableauHeader{Apport pour la preuve}
}
Convention de Malabo (UA, 2014) & Cadre continental~: infractions, données, coopération \\
Convention de Budapest (2001) & Preuve électronique, perquisition, entraide internationale \\
\end{tableaucours}
\vspace{-0.8em}
\begin{center}{\small\itshape Tableau~4.1. Deux instruments supranationaux~\cite{malabo2014,budapest2001}}\end{center}

% ============================================================================
\section{Le Juge et l'Expert}
\label{sec:c04-juge}
% ============================================================================

Le juge apprécie la preuve~; l'expert l'éclaire. L'expert ne tranche pas la
culpabilité~: il établit des faits techniques et explique leur portée. La
frontière est nette, et la respecter protège la valeur de l'expertise.

\begin{boxscenario}{Opération TSANAGA. Le titre d'abord}
\textbf{Constatation.} Avant toute opération, Me~TSALA vérifie que la
réquisition du parquet couvre bien les supports saisis et le périmètre de
l'analyse.

\medskip
\textbf{Analyse} (confiance élevée). Sans ce contrôle, une analyse hors
périmètre aurait pu être écartée. Le droit n'est pas un obstacle à l'enquête~:
c'est ce qui en garantit la valeur.
\end{boxscenario}

\begin{boxmetier}{Juriste TIC}
\textbf{Sécuriser l'enquête juridiquement.} Le juriste vérifie le titre, le
périmètre, le respect des données, et prépare la recevabilité. Son travail,
en amont, évite que des mois d'enquête ne s'effondrent devant le juge.
\end{boxmetier}

\separateur

\begin{boxresume}
\textbf{Ce qu'il faut retenir de ce chapitre~:}
\begin{itemize}
    \item \textbf{Loi 2010/012}~: accès frauduleux, preuve électronique.
    \item \textbf{Loi 2024/017}~: protection des données personnelles.
    \item \textbf{Malabo et Budapest} encadrent preuve et coopération.
    \item \textbf{Le juge apprécie}, \textbf{l'expert éclaire}~; il ne tranche
          pas la culpabilité.
    \item \textbf{Le titre et le périmètre} se vérifient avant d'agir.
\end{itemize}
\end{boxresume}

\begin{boxinfo}
\textbf{Points de contrôle de l'élève.} Je sais~: citer les textes
camerounais~; situer Malabo et Budapest~; distinguer les rôles du juge et de
l'expert~; énoncer les règles à respecter.
\end{boxinfo}

\bigskip

\begin{boxexo}[title={\ding{45}\quad Exercice~4.1. Le texte \niveaubase}]
Quel texte camerounais encadre la preuve électronique~?
\end{boxexo}

\begin{boxexo}[title={\ding{45}\quad Exercice~4.2. Données \niveaubase}]
Quel texte protège les données personnelles~?
\end{boxexo}

\begin{boxexo}[title={\ding{45}\quad Exercice~4.3. Juge ou expert \niveaubase}]
Qui apprécie la preuve~? Qui l'éclaire~?
\end{boxexo}

\begin{boxexo}[title={\ding{45}\quad Exercice~4.4. Sans titre \niveaumoyen}]
Que risque une preuve obtenue sans titre~?
\end{boxexo}

\begin{boxexo}[title={\ding{45}\quad Exercice~4.5. Périmètre \niveaumoyen}]
Pourquoi vérifier que la réquisition couvre le périmètre d'analyse~?
\end{boxexo}

\begin{boxexo}[title={\ding{45}\quad Exercice~4.6. Budapest \niveaumoyen}]
Qu'apporte la Convention de Budapest à une enquête transfrontalière~?
\end{boxexo}

\begin{boxexo}[title={\ding{45}\quad Exercice~4.7. Tiers \niveaumoyen}]
Comment protéger les tiers dont les données figurent sur un support saisi~?
\end{boxexo}

\begin{boxexo}[title={\ding{45}\quad Exercice~4.8. La frontière \niveaumoyen}]
Pourquoi l'expert ne doit-il pas conclure à la culpabilité~?
\end{boxexo}

\begin{boxexo}[title={\ding{45}\quad Exercice~4.9. Note juridique \niveauexpert}]
Rédigez la note qui sécurise juridiquement l'opération TSANAGA.
\end{boxexo}

\begin{boxexo}[title={\ding{45}\quad Exercice~4.10. Comparer \niveauexpert}]
Comparez le cadre camerounais à la Convention de Malabo sur la preuve.
\end{boxexo}

\bigskip

\begin{boxinfo}
\textbf{Pour aller plus loin.}
\begin{itemize}
    \item Lois 2010/012~\cite{loi2010012} et 2024/017~\cite{loi2024017}~;
          Malabo~\cite{malabo2014}~; Budapest~\cite{budapest2001}.
\end{itemize}
\end{boxinfo}

\begin{boxinfo}
\textbf{Transition.} Le droit exige une continuité de la détention~: le
Chapitre~\ref{chap:custody} traite de la chaîne de possession.
\end{boxinfo}
